\section{Основные принципы автоматизации процесса работы наземных транспортно-технологических средств}\label{sec:ch1/sec2}

В общем виде автоматизируются процессы перемещения и выполнения рабочих операций (сбор, сортировка и другие). Это подразумевает автоматизацию следующих операций:

\begin{itemize}
    \item получение и обработка данных с датчиков;
    \item их объединение и согласование;
    \item сбор и обработка изображений;
    \item определение дорожных условий, препятствий и расстояния до них;
    \item позиционирование и определение текущего состояния;
    \item реализация автоматического управления скоростью, траекторией (курсом) движения, автоматической реакцией на окружающие объекты;
    \item принятие решений на управляющие действия;
    \item управление исполнительными устройствами;
    \item сбор и передача телеметрии для последующего анализа.
\end{itemize}

\begin{figure}[ht]
	\centerfloat{
		\includegraphics[scale=0.5]{drawio/arch_tech.drawio.png}
	}
	\caption{Общая архитектура технологий автоматизации \linebreak рабочих~процессов}\label{fig:arch_tech}
\end{figure}

Архитектура технологии в целом представлена на рисунке~\cref{fig:arch_tech} и состоит из взаимосвязанных частей (подсистем):

\begin{itemize}
	\item сбора данных (датчиков и сенсоров, восприятия) (см.\:раздел\:\cref{subsec:ch1/sec2/sub1});
	\item позиционирования и локализации (см.\:\cref{subsec:ch1/sec2/sub2});
	\item исполнительной подсистемы (см.\:\cref{subsec:ch1/sec2/sub3});
	\item управления и контроля (см.\:\cref{subsec:ch1/sec2/sub4}).
\end{itemize}

С помощью различных датчиков РНТТМ обеспечивается сбор данных о состоянии окружающей среды. Датчики можно классифицировать на датчики внешней информации и внутренней \cite{Vlasov}.

Датчики внутренней информации РНТТМ представляют собой в основном преобразователи механических параметров в электрические сигналы. Например датчики обратной связи по положению, скорости и ускорения.

\begin{figure}[ht]
    \centering{
        \includegraphics[scale=0.45]{drawio/datchiki.drawio.png}
    }
    \caption{Классификация датчиков внешней информации РНТТМ}\label{fig:datchiki}
\end{figure}

Датчики внешней информации РНТТМ обеспечивают сбор информации об окружающем мире с различными целями. Датчики можно разделить на две большие группы (рис.\ \cref{fig:datchiki}): дистанционные и контактные. В свою очередь дистанционные датчики можно разделить на локационные и визуальные.

Локационные системы РНТТМ можно условно разделить на два класса: дальней и~ближней локации. Первые могут быть построены с использованием ультразвуковых, лазерных и светолокационных датчиков.

Визуальные системы (технического, машинного зрения) РНТТМ относятся к числу наиболее сложных, но и наиболее универсальных средств очувствления РНТТМ (рис.\ \cref{fig:ypiter30},\ \cref{fig:robot1},\ \cref{fig:robot2}).

\begin{figure}[ht]
	\centering{
		\hfill
		\subcaptionbox{Робота\label{fig:robot1}}{%
			\includegraphics[width=0.4\linewidth]{views/robot1.jpg}}
		\hfill
		\subcaptionbox{Кобота\label{fig:robot2}}{%
			\includegraphics[width=0.4\linewidth]{views/robot2.jpg}}
		\hfill
	}
	\caption{Примеры наборов датчиков}\label{fig:robot1_2}
\end{figure}

Основными используемыми средствами машинного видения РНТТМ~\cite{confbib1} являются:

Лазерные дальномеры (лидары) РНТТМ, которые измеряют расстояние до объектов. Результат работы лидара -- объемная сцена из облака точек с геометрией объектов. Это средство выдает наибольшую точность, лучше всех работают на ровных поверхностях и почти не засвечиваются солнцем. Недостатки лидаров -- это высокое энергопотребление, относительно низкая частота кадров, бегущий затвор и необходимость компенсировать его при обработке, а также работающие рядом лидары создают друг-другу помехи, которые не так просто компенсировать.

Радары РНТТМ для определения расстояния до объекта, его скорости и~местораспо-ложения. Радары меньше зависят от погоды и стоят намного дешевле лидаров. Главным отрицательным качеством радаров является то, что плохо обнаруживаются неметаллические объекты, в том числе пешеходы.

Инфракрасные датчики РНТТМ для реагирования на фоновое инфракрасное излучение. Прибор регистрирует любое тепловое излучение.

Time-of-flight-камеры (ToF) РНТТМ -- видеокамеры, формирующие дальностное изображение. Используются для создания изображений, которые в качестве пикселей содержат оценки расстояний от экрана до конкретных точек наблюдения.

\nomenclature{ToF}{Time-of-flight-камеры\nomrefpage}

Стерео-камеры РНТТМ, построенные на имитации бинокулярного зрения человека с возможностью захватывать трехмерные изображения.

Пленоптические камеры РНТТМ, фиксирующие векторное поле световых лучей (световое поле). На основе картины светового поля может быть воссоздана наиболее полная информация об изображении, пригодная для решения различных задач компьютерной графики.

Обычные камеры РНТТМ с последующим анализом видео-потока с применением алгоритмов машинного обучения.

Применение какого-либо только одного средства недостаточно, так как получение комплексного видения окружающего мира роботами имеет сложности в виде:

\begin{itemize}[beginpenalty=10000, midpenalty=10000, wide=\parindent, leftmargin=0pt]
    \item темного времени суток;
    \item засветов от солнца и источников освещения;
    \item инфракрасного излучение солнечного света;
    \item осадков и грязи;
    \item множественности и вариативности подвижных окружающих объектов.
\end{itemize}

В связи с этим, средства машинного видения в РНТТМ применяются комплексно, для компенсации перечисленных недостатков. Сравнительный анализ приведен в~таблице~\ref{tab:TableViews}. Согласно требованиям к подсистеме распознавания ГОСТ Р 71615-2024 (п. 6.4.10) \cite{gostR71615}, для обеспечения функциональной безопасности ВАЗДСМ необходимо комбинирование указанных средств технического зрения (лидары, радары, камеры) для исключения ложных срабатываний в условиях запыленности и осадков.

\begin{table}[htbp]
	\centering
	\caption{Сравнительный анализ средств машинного видения РНТТМ}\label{tab:TableViews}
	\resizebox{\textwidth}{!}{%
		\begin{threeparttable}
			\begin{tabular}{|l|c|c|c|c|c|}
				\hline
				& \textbf{ToF} & \textbf{Стерео} & \textbf{Пленопт.} & \textbf{Лидары} & \textbf{Радары} \\ \hline
				Разрешение                 & \textit{Слабо}  & Средне & \textbf{Отлично (сцена)} & Средне (сцена) & \textit{Слабо}  \\ \hline
				Точность                  & Средне & Средне & \textit{Слабо} & \textit{Слабо} & \textbf{Высоко} \\ \hline
				Сложность программ        & \textbf{Слабо} & Средне & \textit{Высоко} & \textit{Высоко} & \textbf{Слабо} \\ \hline
				Работа в реальном времени & \textbf{Высоко} & Средне & Средне & Средне & \textit{Слабо} \\ \hline
				Стоимость                 & Средне & Средне & \textbf{Слабо} & \textit{Высоко} & Средне \\ \hline
				Произв. при слабом освещении & \textbf{Хорошо} & \textbf{Хорошо} & \textit{Плохо} & \textit{Плохо} & \textbf{Хорошо} \\ \hline
				Работа на открытом воздухе & \textit{Плохо} & \textit{Плохо} & \textbf{Хорошо} & \textbf{Хорошо} & \textbf{Хорошо} \\ \hline
			\end{tabular}
		\end{threeparttable}
	}
\end{table}

Можно сделать вывод, что по разрешению лидируют пленоптические камеры, но все очень сильно зависит от состояния окружающей среды (сцены съемки). По точности -- лидары в РНТТМ. По сложности обработки -- «непосредственно» получают глубину только ToF и~лидары. По частоте кадров конструктивно хороши ToF-камеры. При работе на открытом пространстве плохо работают ToF и~камеры структурированного света. % Сбор данных

Локализация и позиционирование -- ключевые элементы системы навигации и управления РНТТМ, позволяющий определить их положение и~движение в пространстве. Для определения положения, ориентации, скорости и ускорения РНТТМ могут использоваться различные данные. Как уже было отмечено на рисунке~\cref{fig:arch_tech}, основные компоненты включают:

\begin{itemize}
	\item Global Navigation Satellite System (GNSS)	-- система спутниковой навигации, которая определяет местоположение на поверхности Земли. Она объединяет данные от нескольких спутниковых систем, таких как GPS (США), ГЛОНАСС (РФ), Galileo (EC) и BeiDou (КНР). Предоставляет данные о местоположении, но может быть неточным в условиях плохой видимости спутников, например, в городах или тоннелях;

	\nomenclature{GNSS}{Global Navigation Satellite System\nomrefpage}
	\nomenclature{США}{Соединённые Штаты Америки\nomrefpage}
	\nomenclature{РФ}{Российская Федерация\nomrefpage}
	\nomenclature{ЕС}{Европейский Союз\nomrefpage}
	\nomenclature{КНР}{Китайская Народная Республика\nomrefpage}

 	\item инерциальный измерительный модуль -- устройство, состоящее из акселерометра и гироскопа. Оно измеряет ускорение и угловую скорость, что помогает отслеживать движение робота. Его показания могут искажаться при вибрациях или наклонах;

	\item датчики вращения колес — высокоточные устройства, установленные на колесах. Они фиксируют повороты колес, что позволяет точно измерять пройденное расстояние и скорость. Точность снижается на скользких или неровных поверхностях;

	\item высококачественная картография -- карты высокой детализации, содержащие трехмерные данные об окружающей среде. Они состоят из лидарного и векторного слоев:

	\begin{enumerate}[left=0pt]
		\item \textit{лидарный слой}:
		трехмерная модель, созданная на основе данных лидара. Она формируется автоматически: РНТТМ проезжает по местности, собирает данные, которые затем обрабатываются для создания детальной карты. Динамические объекты, такие как машины или пешеходы, должны удалятся из данных, чтобы карта оставалась стабильной. Для этого может применяется алгоритм Normal Distribution Transform, который сравнивает текущие данные лидара с картой;

		\item \textit{векторный слой}:
		часть высококачественной картографии, содержащая информацию о дорожной инфраструктуре. Этот слой включает данные о пешеходных переходах, тротуарах, велодорожках и других зонах. Также в нем хранится информация о светофорах и их взаимосвязи, что помогает роботу планировать маршрут.
	\end{enumerate}
\end{itemize}

В направлении объединения этих данных имеются известные наработки, например метод Extended Kalman, который соединяет данные от всех сенсоров для получения точного результата. Он учитывает погрешности каждого сенсора и корректирует данные, чтобы минимизировать ошибки. Однако для его работы требуется правильная настройка, иначе могут возникать неточности.

Таким образом, локализация РНТТМ достигается за счет комбинации данных от различных сенсоров и алгоритмов, что позволяет точно определять его положение и движение в пространстве. % Позиционирование и локализация

\subsection{Исполнительная подсистема}\label{subsec:ch1/sec2/sub3}

Любая машина, в том числе и НТТС, состоит из отдельных частей: силового (приводы) и рабочего оборудования, передаточных устройств, системы управления и других.

Исполнительная подсистема (ИcП) является ключевым элементом автоматизации НТТС, обеспечивая непосредственное выполнение управляющих команд, сформированных системой управления. Основная задача ИcП заключается в преобразовании электрических сигналов в механические действия, необходимые для перемещения и выполнения рабочих операций НТТС. Объектами роботизации (автоматизации) в строительстве являются технологические операции, выполняемые при возведении зданий и сооружений. Например процессы земляные, монтажные или отделочные.

\nomenclature{ИcП}{исполнительная подсистема\nomrefpage}

В состав ИcП входят следующие основные компоненты \cite{Vilman, Evtukov}:

\begin{itemize}
	\item \textit{приводы}:
	устройства, преобразующие электрическую энергию в механическое движение. В зависимости от типа НТТС и решаемых задач могут использоваться электродвигатели, гидравлические или пневматические приводы. Электродвигатели обеспечивают высокую точность и быстродействие, в то время как гидравлические и пневматические системы применяются для задач, требующих значительных усилий, например, при управлении тяжелыми механизмами или подъемными устройствами;

	\item \textit{рабочие органы}:
	представляют собой рабочую часть манипулятора, служат для непосредственного выполнения технических операций или вспомогательных переходов. Включают в себя механизмы, отвечающие за выполнение специфических операций, таких как подъем, перемещение или обработка грузов. Например, в случае с автономными погрузчиками или экскаваторами это могут быть гидравлические манипуляторы, ковши или другие рабочие инструменты ;

	\item \textit{тормозные системы}:
	обеспечивают управление скоростью и остановку НТТС. В автоматизированных системах тормозные механизмы могут быть интегрированы с датчиками обратной связи, что позволяет точно контролировать процесс торможения в зависимости от дорожных условий и текущей скорости;

	\item \textit{рулевые механизмы}:
	отвечают за изменение направления движения НТТС. В автоматизированных системах рулевое управление реализуется с помощью электромеханических или гидравлических усилителей, которые обеспечивают точное выполнение команд по изменению траектории движения;

	\item \textit{системы обратной связи}:
	обеспечивают контроль за состоянием исполнительных устройств и передачу данных о выполнении команд в систему управления. Это позволяет корректировать управляющие воздействия в~реальном времени, минимизируя ошибки и повышая точность выполнения операций.
\end{itemize}

Особенностью ИcП НТТС является необходимость обеспечения высокой надежности и отказоустойчивости, что особенно важно в условиях эксплуатации, связанных с повышенными нагрузками и вибрациями. Для этого применяются резервирование критически важных компонентов, использование защищенных от внешних воздействий конструкций, а также интеграция систем диагностики, способных своевременно выявлять и устранять неисправности.

Кроме того, ИcП должна быть адаптирована к специфике работы НТТС, включая возможность управления сложными конструкциями, такими как полуприцепы или сменное рабочее оборудование (рис.\:\cref{fig:manipul}). Это требует разработки специализированных алгоритмов управления, учитывающих изменение динамических характеристик и массово-инерционных параметров системы \cite{PevznerUgaiSulemanov2010}.

\begin{figure}[ht]
	\centering{
		\includegraphics[scale=0.5]{views/manipulator.png}
	}
	\caption{Пример полуприцепа и рабочего оборудования}\label{fig:manipul}
\end{figure}

Одновременно с этим, являясь ключевым элементом автоматизации НТТС, ИcП может оказывать значительное влияние на другие подсистемы, как положительное, так и отрицательное. Это влияние может проявляться в виде ограничений, накладываемых на работу других подсистем, или в виде необходимости адаптации их функциональности для обеспечения согласованной работы всей системы в целом. Рассмотрим основные аспекты этого влияния.

ИcП может создавать дополнительные шумы и вибрации, которые могут негативно сказаться на точности работы датчиков, особенно чувствительных к механическим воздействиям. Например, вибрации от работы двигателей или гидравлических систем могут искажать данные акселерометров и гироскопов, что затрудняет точное определение положения и ориентации НТТС.

Электромагнитные помехи от мощных электродвигателей или силовых кабелей могут влиять на работу радаров, лидаров и других электронных датчиков \cite{confbib1}.

Для минимизации этих эффектов требуется тщательное проектирование системы, включая экранирование датчиков, использование фильтров для подавления шумов и вибраций, а также алгоритмов компенсации помех.

ИcП может влиять на точность позиционирования и локализации, особенно в динамических условиях: изменение массы и центра тяжести при работе с грузами или сменным оборудованием (например, полуприцепами) может привести к смещению расчетных параметров движения, что требует корректировки алгоритмов локализации.

Задержки в выполнении команд (например, из-за инерции тяжелых механизмов) могут привести к рассинхронизации между расчетной и фактической траекторией движения, что особенно критично при работе в условиях ограниченного пространства.

Для устранения этих проблем необходимо учитывать динамические характеристики ИcП при разработке алгоритмов управления и локализации, а также использовать системы обратной связи для оперативной корректировки.

ИcП может накладывать ограничения на работу подсистемы управления, особенно в части быстродействия -- механические системы (например, гидравлические приводы) могут иметь ограниченную скорость реакции, что требует адаптации алгоритмов управления для учета этих задержек. Энергопотребления -- мощные исполнительные устройства (например, двигатели или гидравлические насосы) могут потреблять значительное количество энергии, что требует оптимизации работы других подсистем для обеспечения баланса энергопотребления. Надежности -- отказ исполнительной подсистемы (например, заклинивание рулевого механизма или потеря управления тормозами) может привести к критическим последствиям, что требует разработки резервных контуров управления и систем диагностики.

ИcП может создавать дополнительные ограничения при работе в сложных условиях:

\begin{itemize}
	\item \textit{по маневренности}: крупные габариты и сложная динамика НТТС (особенно при наличии полуприцепов или сменного оборудования) могут затруднять движение в узких пространствах или на пересеченной местности;

	\item \textit{по скорости}: необходимость обеспечения безопасности и устойчивости НТТС может требовать снижения скорости движения, что влияет на общую производительность системы.
\end{itemize}

Для снижения негативного влияния ИcП на другие подсистемы могут быть применены следующие меры:

\begin{itemize}
	\item \textit{интеграция систем обратной связи}: для оперативного учета изменений в работе исполнительных устройств и корректировки управляющих команд;

	\item \textit{оптимизация алгоритмов управления}: с учетом динамических характеристик ИcП и возможных задержек;

	\item \textit{резервирование}: для повышения надежности и отказоустойчивости системы;

	\item \textit{адаптация под условия эксплуатации}: например, использование специализированных алгоритмов для работы с полуприцепами или в условиях пересеченной местности.
\end{itemize}

Таким образом, ИcП, хотя и является важным элементом автоматизации НТТС, может накладывать значительные ограничения на работу других подсистем. Однако при грамотном проектировании и интеграции этих ограничений можно минимизировать, обеспечивая согласованную и эффективную работу всей системы в целом. % исполнительной подсистемы

В общем случае различают НТТС с механизированным и автоматизированным управлением, но в \cite[с.\ 39]{Evtukov} рассмотрено шире (рис.\ \cref{fig:control_drawio}).

\begin{figure}[ht]
	\centering{
		\includegraphics[scale=0.5]{drawio/control.drawio.png}
	}
	\caption{Классификация систем управления и контроля НТТС}\label{fig:control_drawio}
\end{figure}

Все реализации систем с автоматизированным управлением используют такие компоненты, как вычислительные системы, сети связи (локальные и глобальные) и программное обеспечение (общее и специальное).

\subsubsection{Вычислительные системы}\label{subsec:ch1/sec2/sub4/sub1}
\subsubsection{Вычислительные системы}\label{subsec:ch1/sec2/sub4/sub1}

\fixme{
	пробить по источникам информации, добавить и плюс картинки для доклада
}

В автоматизированной системе управления НТТС могут использоваться следующие элементы вычислительных систем \cite{Tanenbaum, MunisterSoc, DigitalDesign} (рис.\ \cref{fig:sensors}):

\begin{itemize}[leftmargin=*]
	\item \textit{центральные процессоры (CPU)}:
	используются для выполнения общих задач, таких как управление операционной системой, обработка данных от датчиков и выполнение алгоритмов (процессоры Intel, AMD);

	\item \textit{графические процессоры (GPU)}:
	используются для задач машинного обучения и компьютерного зрения, таких как распознавание объектов, классификация изображений и обработка данных с камер (NVIDIA Drive AGX, NVIDIA Tesla, A100);

	\item \textit{тензорные процессоры (TPU)}:
	это специализированные процессоры для ускорения вычислений в нейронных сетях (Google TPU);

	\item \textit{системы на кристалле (SoC)}:
	более компактные решения, объединяющие CPU, GPU, TPU и другие компоненты на одном чипе (NVIDIA Xavier, NVIDIA Orin, Qualcomm Snapdragon Ride);

	\item \textit{программируемые логические интегральные схемы (FPGA)}:
	необходимы для задач, требующих высокой производительности и низкого энергопотребления, таких как обработка сигналов с лидаров и радаров (Xilinx Zynq, Intel FPGA);

	\item \textit{специализированные процессоры для автономных систем}:
	специальные для беспилотных транспортных средств, которые обеспечивают высокую производительность и энергоэффективность (NVIDIA Drive PX, Mobileye EyeQ);

	\item \textit{датчики и их интерфейсы}:
	лидары, радары, камеры, ультразвуковые датчики и GPS-модули подключаются к вычислительным системам через специализированные интерфейсы (CAN, Ethernet, PCIe);

	\item \textit{оперативная память}:
	высокоскоростная память для временного хранения данных, необходимых для обработки в реальном времени (DDR4, GDDR6);

	\item \textit{накопители данных}:
	используются для хранения больших объемов данных, таких как карты, модели машинного обучения и логи (NVMe SSD);

	\item \textit{коммуникационные модули}:
	обеспечение связи между компонентами системы, а также с внешними серверами (Ethernet, 5G-модули);

	\item \textit{микроконтроллеры}:
	для управления низкоуровневыми системами, такими как двигатели, тормоза и рулевое управление (ARM Cortex-M, STM32).
\end{itemize}

\nomenclature{CPU}{Central Processing Unit, центральный процессор\nomrefpage}
\nomenclature{GPU}{Graphics Processing Unit, графический процессор\nomrefpage}
\nomenclature{TPU}{Tensor Processing Unit, тензорный процессор\nomrefpage}
\nomenclature{SoC}{System on a Chip, электронная схема на одном кристалле\nomrefpage}
\nomenclature{FPGA}{field-programmable gate array, программируемая пользователем вентильная матрица\nomrefpage}

\begin{figure}[ht]
	\centering{
		\includegraphics[scale=0.6]{views/sensors.png}
	}
	\caption{Вариант схемы взаимодействия компонент системы автоматизированного управления в НТТС}\label{fig:sensors}
\end{figure}

\subsubsection{Сети связи}\label{subsec:ch1/sec2/sub4/sub2}
Сети связи -- отдельная подсистема автоматизированной системы управления (АСУ) НТТС и состоит из внутренней сети и связи с внешними источниками (рис.\ \cref{fig:ntts_link}). Они играют ключевую роль и обеспечивают взаимодействие между внутренними компонентами системы и внешними источниками данных соответственно. На рис.\ \ref{fig:sensors} представлена схема взаимодействия компонентов системы автоматизированного управления НТТС, где сети связи играют ключевую роль. Рассмотрим физический и логический уровни сетей связи отдельно, чтобы подчеркнуть их особенности и важность для НТТС \cite{Martin}.

\nomenclature{АСУ}{автоматизированная система управления\nomrefpage}

\begin{figure}[ht]
	\centerfloat{
		\includegraphics[width=14cm]{drawio/ntts_link.drawio.png}
	}
	\caption{Классификация сетей связи АСУ НТТС}\label{fig:ntts_link}
\end{figure}

Физический уровень определяет способы передачи данных через физические среды, такие как радиоволны, оптические волокна или проводные соединения. На этом уровне реализуются механизмы контроля целостности передаваемых данных в конкретном фрейме (кадре). Но не реализуются механизмы адресования, гарантий доставки данных и т.\ д.

Для внутренней сети связи используются исключительно проводные соединения в виду их надежности и стабильности работы, такие как:
\begin{itemize}
	\item \textit{универсальная последовательная шина (USB\ \footnote{\  Последовательный интерфейс для подключения периферийных устройств к вычислительной технике.})}:
	стандарт для подключения периферийных устройств (камер, датчиков и т.\ д.) к вычислительной технике (см.\ раздел\ \ref{subsec:ch1/sec2/sub1});

	\item \textit{сеть контроллеров (CAN-шина\ \footnote{\ Стандарт промышленной сети, ориентированный, прежде всего, на объединение в единую сеть различных исполнительных устройств и датчиков. Режим передачи — последовательный, широковещательный, пакетный.})}:
	является стандартом для автомобильной автоматики с 1980х годов и широко апробирован в различных вариантах использования;

	\item \textit{проводной Ethernet}:
	самый распространенный и надежный протокол передачи данных. Работает на канальном и физическом уровне модели взаимодействия открытых систем (OSI) \cite{Olifers, TCP_IP, ComputerNetwork} (рис.\ \cref{fig:osi_tcp}).
\end{itemize}

\nomenclature{USB}{Universal Serial Bus, универсальная последовательная шина\nomrefpage}
\nomenclature{OSI}{Open Systems Interconnection model, сетевая модель стека сетевых протоколов OSI/ISO\nomrefpage}
\nomenclature{CAN}{Controller Area Network, сеть контроллеров\nomrefpage}

\begin{figure}[ht]
	\centerfloat{
		\includegraphics[scale=0.3]{views/osi.jpg}
	}
	\caption{Сравнение моделей протоколов передачи данных}\label{fig:osi_tcp}
\end{figure}

Так как НТТС -- это мобильные объекты, то наиболее актуальными для связи с внешними источниками (в том числе друг с другом) являются беспроводные технологии, такие как сети мобильной связи (типа
4G\ \footnote{\ \ Поколение мобильной связи с повышенными требованиями. К четвертому поколению принято относить технологии, позволяющие осуществлять передачу данных со скоростью до 100 Мбит/с подвижным (с высокой мобильностью) и до 1 Гбит/с — стационарным абонентам (с низкой мобильностью).}, 5G\ \footnote{\ \ Поколение мобильной связи, действующее на основе стандартов телекоммуникаций, следующих за существующими стандартами 4G.} и др. (рис.\ \cref{fig:cellular})), сети Wi-Fi\ \footnote{\ \ Технология беспроводной локальной сети с устройствами на основе стандартов IEEE 802.11.}, спутниковой связи (например серии Ямал (рис.\ \cref{fig:yamal}) или Starlink (рис.\ \cref{fig:starlink})) и радиосвязь УКВ.

\nomenclature{4G}{fourth generation, четвертое поколение мобильной связи\nomrefpage}
\nomenclature{5G}{fifth generation, пятое поколение мобильной связи\nomrefpage}
\nomenclature{УКВ}{ультракороткие волны диапазона радиоволн\nomrefpage}

Области применения этих технологий в контексте НТТС:
\begin{itemize}
	\item сети мобильной связи обеспечивают высокую пропускную способность и низкую задержку, что критически важно для управления НТТС в~режиме реального времени \cite{6G};
	\item спутниковая связь способна обеспечивать коммуникацию в удаленных и малонаселенных районах, где отсутствует покрытие наземных сетей мобильной связи \cite{SystemComm}. При этом для движущихся объектов, таких как НТТС, необходимо использовать специализированные спутниковые системы, обеспечивающие стабильное соединение даже в условиях высокой мобильности;
	\item Wi-Fi имеет ограниченный радиус применения \cite{Olifers}, можно использовать в помещениях, для локальной связи друг с другом. Так же обеспечивают высокую пропускную способность и низкую задержку;
	\item УКВ целесообразно рассматривать как вид резервной связи, т.\ к. этот вид связи подразумевает низкую пропускную способность, но обеспечит передачу данных при отсутствии других видов связи.
\end{itemize}

Качественное сравнение представлено в табл.\ \ref{tab:PhysicalLevel}.

\begin{figure}[ht]
	\centerfloat{
		\includegraphics[scale=0.2]{views/cell.png}
	}
	\caption{Поколения сетей мобильной связи}\label{fig:cellular}
\end{figure}

\begin{figure}[ht]
	\centerfloat{
		\includegraphics[width=14cm]{views/yamal.jpg}
	}
	\caption{Зоны покрытия спутников сети Ямал}\label{fig:yamal}
\end{figure}

\begin{figure}[ht]
	\centerfloat{
		\includegraphics[width=14cm]{views/starlink.png}
	}
	\caption{Зоны работы спутников сети Starlink}\label{fig:starlink}
\end{figure}

\begin{table}[htbp]
	\centering
	\caption{Сравнение технологий беспроводной связи}\label{tab:PhysicalLevel}
	\resizebox{\textwidth}{!}{%
		\begin{threeparttable}
			\begin{tabular}{|c|c|c|c|c|c|}
				\hline
					\textbf{Технология} & \textbf{Покрытие} & \textbf{Надежность} & \textbf{Скорость} & \textbf{Задержка} & \textbf{Применение} \\ \hline
				\makecell{4G/5G} & \makecell{Городские\\районы} & Высокая & \makecell{1-10 Гбит/с} & \makecell{1-50 мс} & \makecell{Управление\\в реальном времени} \\ \hline
				\makecell{Спутниковая} & Глобальное & Средняя & \makecell{100-500 Мбит/с} & \makecell{20-800 мс} & \makecell{Удаленные\\районы} \\ \hline
				Wi-Fi & \makecell{Локальное} & Высокая & \makecell{3.5-9.6 Гбит/с} & \makecell{1-20 мс} & \makecell{Локальная связь} \\ \hline
				УКВ & \makecell{До 50 км} & Низкая & \makecell{до 10 Мбит/с} & \makecell{100-500 мс} & \makecell{Резервная связь} \\ \hline
			\end{tabular}
		\end{threeparttable}
	}
\end{table}

Логический уровень определяет протоколы и методы передачи данных, которые обеспечивают надежность, целостность и эффективность связи. Он основан и использует физический, но при этом абстрагирован, т.\ к. данные логического уровня инкапсулируются в кадры (фреймы) физического (рис.\ \cref{fig:tcp_udp}). На логическом уровне реализуются адресование (глобальное и локальное) и~ механизмы логического соединения между прикладными процессами, выполняющимися на разных объектах. При этом логическое соединение с точки зрения приложений выглядит как канал, непосредственно соединяющий процессы.

\begin{figure}[ht]
	\centerfloat{
		\includegraphics[width=8cm]{views/tcp_udp.png}
	}
	\caption{Инкапсуляция данных}\label{fig:tcp_udp}
\end{figure}

В контексте НТТС следует рассматривать:
\begin{itemize}
	\item протоколы адресования, которые реализуют механизмы поиска путей пакетов в сети:
		\begin{enumerate}
			\item межсетевые протоколы IP\ \footnote{\ \ Маршрутизируемый протокол сетевого уровня стека TCP/IP.},
			OSPF\ \footnote{\ \ Основан на технологии отслеживания состояния канала (link-state technology) и использующий для нахождения кратчайшего пути алгоритм Дейкстры.},
			BGP\ \footnote{\ \ Протокол динамической маршрутизации, он же eBGP, от exterior. Относится к классу протоколов маршрутизации внешнего шлюза.},
			IS-IS\ \footnote{\ \ Протокол внутренних шлюзов (Interial Gateway Protocol), стандартизированный ISO и использующийся в основном в крупных сетях провайдеров услуг.}
			и др. \cite{ComputerNetwork};

			\item единого сетевого (стека) протоколов информационного обмена (ЕСПИО) \cite{ESPIO};

			\nomenclature{IP}{Internet Protocol, протокол интернета\nomrefpage}
			\nomenclature{OSPF}{Open Shortest Path First, протокол динамической маршрутизации\nomrefpage}
			\nomenclature{BGP}{Border Gateway Protocol, протокол граничного шлюза\nomrefpage}
			\nomenclature{IS-IS}{протокол маршрутизации промежуточных систем\nomrefpage}
			\nomenclature{ЕСПИО}{единый стек протоколов информационного обмена\nomrefpage}
		\end{enumerate}

	\item транспортные протоколы:
		\begin{enumerate}
			\item MQTT\ \footnote{\ \ Упрощенный сетевой протокол, работающий поверх TCP/IP, ориентированный на обмен сообщениями между устройствами по принципу <<издатель -- подписчик>>.}:
			легковесный протокол для IoT\ \footnote{\ \ Концепция сети передачи данных между физическими объектами, оснащенными встроенными средствами и технологиями для взаимодействия друг с другом или с внешней средой.} с поддержкой QoS \cite{Munister};

			\item LoRaWAN\ \footnote{\ \ Беспроводная технология передачи небольших по объему данных на дальние расстояния, разработанная для распределенных сетей телеметрии, межмашинного взаимодействия и интернета вещей.}:
			энергоэффективный протокол для дальнодействующих сетей \cite{PerryLi};

			\item TCP:
			надежная передача с контролем перегрузок \cite{TCP_IP};

			\item UDP:
			минималистичный протокол без гарантий доставки \cite{ComputerNetwork};

			\item QUIC\ \footnote{\ \ Разработан Google для компенсации недостатков протокола TCP в связи с проблема развертывания новых реализаций TCP на существующих сетях связи. Позволяет мультиплексировать несколько потоков данных между двумя компьютерами, работая по протоколу UDP.}:
			современный протокол с мультиплексированием поверх UDP \cite{Olifers};

			\item MultipathTCP\ \footnote{\ \ Конкурент QUIC. Реализован как часть TCP.}:
			расширение TCP с поддержкой множества путей \cite{ComputerNetwork}.

			\nomenclature{MQTT}{message queuing telemetry transport, протокол передачи очередей сообщений\nomrefpage}
			\nomenclature{IoT}{internet of things, интернет вещей\nomrefpage}
			\nomenclature{QoS}{quality of service, качество обслуживания\nomrefpage}
			\nomenclature{LoRaWAN}{Low-power Wide-area Network, энергоэффективная сеть дальнего радиуса действия\nomrefpage}
			\nomenclature{UDP}{User Datagram Protocol, протокол пользовательских датаграмм\nomrefpage}
			\nomenclature{TCP}{Transmission Control Protocol, протокол управления передачей\nomrefpage}
			\nomenclature{QUIC}{Quick UDP Internet Connections, быстрые UDP подключения по IP\nomrefpage}
		\end{enumerate}
\end{itemize}

НТТС являются потребителями услуг, предоставляемых уровнем протоколов адресования и используют только интерфейс взаимодействия с сетевыми устройствами. Одновременно с этим, деятельность НТТС непрерывно связано с рисками, то очень важно понимать механизмы, обеспечивающие надежную передачу данных транспортными протоколами и возможности их использования:

\begin{itemize}
	\item \textit{контрольная сумма}: используется для обнаружения битовых ошибок в переданном пакете;

	\item \textit{таймер}: отсчет интервала ожидания и указание на его истечение. Последнее означает, что с высокой степенью вероятности пакет или его квитанция потеряны при передаче. В случае, если пакет доставляется с задержкой, но не теряется (преждевременное истечение интервала ожидания), либо происходит потеря квитанции, повторная передача приводит к дублированию пакета на принимающей стороне;

	\item \textit{порядковый номер сегмента}: используется для порядковой нумерации пакетов данных передаваемых от отправителя к получателю. Разрывы в порядковых номерах полученных пакетов позволяют получателю обнаружить потерю пакета. Одинаковые порядковые номера пакетов означают, что пакеты дублируют друг друга;

	\item \textit{подтверждение получения}: генерируется принимающей стороной и указывает передающей стороне на то, что соответствующий пакет или группа пакетов успешно приняты. Обычно подтверждение содержит порядковые номера успешно принятых пакетов. В зависимости от протокола различают индивидуальные и групповые подтверждения;

	\item \textit{отрицательное подтверждение}: используется получателем для сообщения отправителю, что пакет был получен некорректно. Отрицательное подтверждение обычно включает в себя порядковый номер пакета, который не был корректно получен;

	\item \textit{конвейеризация}: ограничивает диапазон порядковых номеров, которые могут использоваться для передачи пакетов. Групповая передача и рукопожатие позволяют значительно увеличить пропускную способность протоколов по сравнению с режимом ожидания подтверждений. Размер окна может быть рассчитан на основе возможностей приема и буферизации принимающей стороны, а также уровня загрузки сети.
\end{itemize}

Качественное сравнение представлено в табл.\ \ref{tab:LogicalLevel}.

 \begin{table}[htbp]
 	\centering
 	\caption{Сравнение протоколов на транспортном уровне}\label{tab:LogicalLevel}
 	\resizebox{\textwidth}{!}{%
 		\begin{threeparttable}
 			\begin{tabular}{|l|c|c|c|c|c|c|c|c|}
 				\hline
 				Протокол & Надежность & \makecell{Эффективность\\ использования\\ канала} & \makecell{Контрольная\\ сумма} & Таймер & \makecell{Подтверждение\\ получения} & \makecell{Отрицательное\\ подтверждение} & Конвейер & Применение \\ \hline
 				MQTT & \makecell{Частичная\\ (QoS 0-2)} & Низкая & Нет & Нет & \makecell{Да\\ (QoS \geq 1)} & Нет & Нет & \makecell{IoT-датчики,\\ телеметрия} \\ \hline
 				LoRaWAN & Средняя & Низкая & Да & Да & Опция & Нет & Нет & \makecell{Удаленные\\ сенсоры} \\ \hline
 				TCP & Высокая & Средняя & Да & Да & Да & Да & Да & \makecell{Веб, передача\\ файлов} \\ \hline
 				UDP & Низкая & \textbf{Макс.} & Опц. & Нет & Нет & Нет & Нет & \makecell{Потоковые данные,\\ VoIP} \\ \hline
 				QUIC & Высокая & Высокая & Да & Да & Да & Да & Да & \makecell{Мобильные\\ приложения} \\ \hline
 				MultipathTCP & \textbf{\makecell{Очень\\ высокая}} & Высокая & Да & Да & Да & Да & Да & \makecell{Критические\\ системы} \\ \hline
 			\end{tabular}
 		\end{threeparttable}
 	}
 \end{table}

\subsubsection{Программное обеспечение}\label{subsec:ch1/sec2/sub4/sub3}
\subsubsection{Программное обеспечение}\label{subsec:ch1/sec2/sub4/sub3}

Программное обеспечение является основой для автоматизации НТТС. Оно позволяет выполнять широкий спектр задач: от управления движением до взаимодействия с внешней средой. Эффективность работы НТТС напрямую зависит от качества используемого ПО, которое делится на несколько категорий (рис.\:\cref{fig:rms_drawio}).

\begin{figure}[ht]
	\centerfloat{
		\includegraphics[width=14cm]{drawio/rms.drawio.png}
	}
	\caption{Программная система управления НТТС}\label{fig:rms_drawio}
\end{figure}

ПО для управления транспортными средствами \cite{LinuxKetov, LinuxWard}. Данная категория включает системы, которые обеспечивают непосредственное управление транспортным средством, включая управление движением, манипуляторами и исполнительными механизмами.

Примеры:

Системы автоматического вождения (например, автопилоты для карьерных самосвалов).
Контроллеры манипуляторов для погрузочно-разгрузочных операций.
Эти системы часто применяются на уровне отдельных транспортных средств и имеют ограниченную функциональность по взаимодействию с другими элементами системы.
ПО для управления процессами
Эти системы отвечают за координацию и управление задачами на уровне всего производственного процесса.
Направления:

Системы планирования маршрутов для транспортных средств в логистических цепочках.
Системы диспетчеризации, которые распределяют задачи между машинами и контролируют их выполнение.
Примером является система управления автопарком (Fleet Management System, FMS), применяемая в логистике и горнодобывающей промышленности.
ПО для анализа и прогнозирования
Данные системы используют машинное обучение и методы анализа данных для прогнозирования и оптимизации работы НТТС.
Примеры:

Прогнозирование состояния транспортных средств для предотвращения поломок (predictive maintenance).
Оптимизация потребления топлива на основе анализа маршрутов и условий работы.
Такие системы часто интегрируются с облачными платформами для доступа к большому объёму данных.
ПО для взаимодействия в многоагентных системах
Данная категория фокусируется на создании координированных действий между несколькими участниками процесса.
Направления:

Системы, обеспечивающие координацию автономных транспортных средств в рамках одного объекта (например, логистического склада).
ПО для взаимодействия автономных транспортных средств с внешней инфраструктурой, включая системы связи V2X (vehicle-to-everything).
Перспективные направления разработки ПО
Интеграция искусственного интеллекта (ИИ)

Использование ИИ для динамического перераспределения задач между участниками.
Реализация систем глубокого обучения для адаптации к меняющимся условиям работы.
Облачные платформы для управления процессами

Централизованное управление данными и задачами через облачные системы.
Интеграция различных типов ПО для создания единого цифрового пространства.
Системы коллективного управления
Перспективными являются системы, которые обеспечивают совместное выполнение задач группой автономных транспортных средств. Например, автономные конвои для логистики или синхронная работа транспортных средств на карьерах.

Универсальные стандарты и открытые платформы

Создание стандартизированных решений, которые позволяют разным системам взаимодействовать друг с другом.
Открытые программные интерфейсы (API) для упрощения интеграции ПО. % управления и контроля